% ===========================================================================
%  08_conclusion.tex -- 结论
% ===========================================================================
\section{结论}\label{sec:concl}

本文从一个很具体的工程目标出发：做出一个\emph{物理上站得住}的黑洞成像模拟器，
让它在普通消费级显卡上实时运行，并且让它的每一个数字都能被独立重算。
围绕这个目标，本文给出了施瓦西时空零测地线的完整推导
（第~\ref{sec:theory} 节）、几何自适应 RK4 积分器的误差分析
（第~\ref{sec:numerics} 节）、单遍片元着色器的 GPU 实现
（第~\ref{sec:impl} 节）、六组逐点验证
（第~\ref{sec:verify} 节）、以及把画面与代价同时变成可复现测量的结果部分
（第~\ref{sec:results} 节）。本节不再引入新的公式或新的实验，
而是把散落在各处的结论收束成三句话，并把它们的\emph{有效范围}
一并说清楚——因为一个物理模拟器的可信度，取决于它对自己残差的诚实程度。

% --------------------------------------------------------------------------- %
\subsection{三句话结论}

\paragraph{第一句：闭式能被反解到机器精度附近，所以「算错了」是会被抓住的。}
本文所有的解析参照量都不是用来说服读者的，而是用来\emph{抓 bug} 的判据。
临界冲击参数从像面反解的相对误差为 $1.0\times10^{-14}$
（第~\ref{sec:v1} 节），这已经是 double 精度下 $\bc=3\sqrt3\,\M$
所能承载的极限；阴影半径的亚像素反解与闭式
$\sin\psi_{\rm sh}=\bc\sqrt{1-\rs/r_0}/r_0$ 相符到 $6.5\times10^{-12}$
（第~\ref{sec:v3} 节，式~\eqref{eq:subpix}）；界面内建的实时测量在
$1440\times820$ 下回读 $144.73$\,px，对解析值 $144.72$\,px 的相对偏差
$+0.009\%$，折算成亚像素残差 $0.013$\,px，与覆盖率量化带
$1/(2\sqrt2\cdot24)=0.0147$\,px 同量级——也就是说，
\emph{观测到的偏差已经落到算法自身的量化地板之下}，
此时任何进一步的「精度提升」都不再是有意义的声明。
近临界轨道的首次积分守恒误差在 $tol=3\times10^{-4}$ 处触底于
$2.6\times10^{-6}$（第~\ref{sec:v2} 节），说明积分器在最容易失稳的
区域也没有出现结构性漂移；有限逃逸半径带来的出射方向误差不超过
$0.021$\,px（第~\ref{sec:v5} 节，式~\eqref{eq:remdef}），
因此把 $r_{\rm esc}$ 取成有限值是安全的。

\paragraph{第二句：代价是线性的、饱和的、可预测的。}
单帧时间满足 $T\approx k\,P+c$，其中 $P$ 为受追踪像素数；
在 $9.6\times10^{4}$ 到 $6.0\times10^{5}$\,px 的动态范围上，
$12$ 个有效测点给出 $k=6.30\times10^{-4}$\,ms/px、
$c=5.2$\,ms，残差 RMS 为 $4.01$\,ms（相对 $2.37\%$，
式~\eqref{eq:costpix}）。这一条\emph{线性律}之所以重要，
是因为它意味着渲染预算可以事先分配：把 \texttt{renderScale}
从 $1.0$ 降到 $0.45$ 会得到 $4.9$ 倍的像素缩减，
而画面的结构（阴影半径以像素计会缩小，但以物理角计不变）
保持不变。与之并行的还有两条标度律：步数在 $tol\gtrsim10^{-2}$
区间的线性响应 $T\propto tol^{-0.912}$（$12$ 个测点，
对数残差 RMS $0.019$，式~\eqref{eq:costtol}），以及
$\rho=\beta/tol$ 参数化下 $25$ 组配置的预算比全部落在
$[0.895,1.089]$ 之内——即\emph{预算参数是饱和的}，
不存在「偷偷变慢」的隐藏通道。需要说明的是，本文的计时是在
一台只有集成显卡（Intel UHD）的机器上、经 ANGLE/D3D11
\emph{硬件}路径测得的；这套管线本身就是 GPU 片元着色器负载，
在配备 NVIDIA 独立显卡的机器上走的是同一个 ANGLE/D3D11 后端、
只是 GPU 算力不同，因此本文给出的绝对时间是\emph{保守上界}而非典型值，
而三条标度律本身与 GPU 型号无关。
在 $tol=0.0075$ 的四组测点上，成本被 $n_{\rm max}=4096$
的步数上限截断（$121.3/487.1/867.4/2326.1$\,ms），
这是设计上的保护而非误差：宁可放弃对极端近临界光线的细分，
也不让单帧无界延长。

\paragraph{第三句：画面上真正出现的东西是可以被\textbf{穷举}的。}
第~\ref{sec:orders} 节对 $300\times169$ 的像面做了逐像素的
交会计数：在 $i=8^\circ$ 时，总 $50700$ 条光线中
$2796$ 条被视界捕获、$17760$ 条与赤道面不相交、
$28502$ 条与盘相交一次、$1642$ 条相交两次，
而相交三次及以上的光线数为 $\mathbf{0}$；在 $i=55^\circ$ 时
同样为 $0$。换言之，在本文的参数域内，
屏幕上只可能出现\emph{一级像与二级像}：
「三四级像」属于近视界区域的相空间，落进阴影以内，
在物理上不可见。这个结论把一个常见的视觉误解
（把阴影边缘的薄环误认为是高阶像）变成了一个可以被数出来的
事实。与之配套的是临界曲线的\emph{对数窄带}：以无量纲距离
$\delta_X$ 逼近临界曲线时，绕行相位的发散斜率在
$\ln\kappa$ 与 $\ln\delta_X$ 两个坐标下分别为
$-1.999946$（残差 RMS $1.95\times10^{-4}$）与
$-0.999944$（$4.37\times10^{-4}$），
推出 $\kappa^2\propto\delta_X$ 的指数为 $2.000059$，
即经典的 $-2$ 与 $-1$ 两个普适斜率
（式~\eqref{eq:slopekappa}--\eqref{eq:expx}、图~\ref{fig:25}）。
从 $\delta_X=10^{-3}$ 到 $10^{-8}$，绕行角只增加
$11.51$\,rad（$1.83$ 圈），却在屏幕尺度上把半径压缩了五个数量级——
这正是近临界成像「圈数有限、梯度无限」的定量表达，
也解释了为什么自适应步长必须以 $u/|\dd u|$ 为标尺
（第~\ref{sec:stepsize} 节）。

% --------------------------------------------------------------------------- %
\subsection{适用范围的明确边界}

上面三句话的每一条都依附于一组前提，把它们写在一起是本文
第~\ref{sec:discuss} 节的主要工作，此处只做提要式的复述：

\begin{itemize}[leftmargin=1.7em,itemsep=2pt,topsep=3pt]
  \item \textbf{时间和自旋。}本文只解静态、球对称、不带电的施瓦西时空。
        Kerr 时空至少改变四类可观测结构（$\bisco$ 位置、阴影形状、
        光子壳层的 Lyapunov 指数、参考系拖曳），其中任何一项都超出
        本文模型的能力（第~\ref{sec:discuss-phys} 节）。
        本文的斜率 $-2$ 与 $-1$ 是施瓦西时空的普适常数，
        不能直接搬到 Kerr。
  \item \textbf{盘与辐射转移。}吸积盘是几何薄的赤道面盘，
        通量取自 Page--Thorne~\cite{pt1974} 或 Shakura--Sunyaev~\cite{ss1973}
        的解析剖面，辐射转移是前向累积的体积积分
        （式~\eqref{eq:front2back}），因此\emph{不含}自吸收以外的
        散射、不含康普顿化、不含盘对射线的自遮挡造成的多普勒遮蔽
        （第~\ref{sec:discuss-rt} 节）。集束律取 $I\propto g^{4}$
        而非 $g^{3}$，两者的差别在第~\ref{sec:discuss-rt} 节给出；
        本文全程使用 $g^{4}$，并在图~\ref{fig:17}(a) 的误差预算中
        把它与离散化、逃逸截断、单精度漂移并列，供读者按需替换。
  \item \textbf{星场。}背景星场只具有\emph{几何}上的可验证性：
        引力透镜对恒星位置的映射可以被解析检验，
        而恒星光度分布是艺术化的，不承担物理结论
        （第~\ref{sec:discuss-stars} 节）。任何基于本模拟器
        恒星亮度的论断都应视为无效。
  \item \textbf{测量不确定度。}第~\ref{sec:discuss-quant} 节给出了
        计时量化与像素反解的不确定度：单点计时在集显机器上会出现
        最高 $24.14\times$ 的组内极差（$36$ 组、每组 $n=7$），
        $27/36$ 组超过 $5\times$，因此本文所有涉及时间的结论
        一律建立在 $n=7$ 的\emph{中位数}之上，而不是单次采样；
        像素反解的量化带为 $1/(2\sqrt2\cdot24)=0.0147$\,px，
        角度分辨率 $\varepsilon_{\rm px}=592.72\,\varepsilon_{\rm rad}$。
        这些数字本身也写进了论文，而不是被一句「测量误差可忽略」掩盖过去。
\end{itemize}

一个顺带的、但值得单独指出的结论来自第~\ref{sec:discuss-lessons} 节：
在写作与验证本文的过程中，共有七处实现缺陷被\emph{物理判据}
而非视觉观察发现——预设系统的状态残留、切换预设后 \texttt{showDisk}
未复位、通量模型开关未进入累积键、\texttt{diskTemp} 在显示模式下
被当作色温使用、单帧成本统计口径不一致、以及抖动种子从未随时间推进。
第七处略有不同：阴影半径探针的亚像素估计器被一处钳位静默退化为
整数扫描（式~\eqref{eq:subpix}），它由跨 $r_0$ 扫描的残差分布形状
而非单点比对抓出，修正后绝对残差上界由 $0.46$\,px 降到 $0.10$\,px。
它们都有一个共同特征：\emph{画面看起来是对的}。
平均亮度、阴影半径、盘的温度渐变都不会因为上述任何一处
而出现肉眼可辨的异常，只有把输出与解析解或与「同一配置应给出同一张图」
这一条件逐点比对时，它们才暴露出来。这七处修正均已同步至开源仓库，
并构成第~\ref{sec:intro} 节贡献列表中的最后一项。

% --------------------------------------------------------------------------- %
\subsection{交付物}

本文所描述的模拟器以\emph{便携式文件夹}的形式交付：
解压后直接双击其中的可执行文件即可运行，无需安装 Node.js、
无需配置 WebGL 环境、无需联网。程序在启动时创建窗口并加载
一篇内嵌的 WebGL2 页面，所有着色器、预设与资源都随包携带；
渲染后端固定走 ANGLE/D3D11 的硬件路径：在本文的测试机上由
Intel UHD 集显执行，在配备 NVIDIA 独立显卡的机器上则由同一后端
直接驱动独显，两条路径逐像素一致（第~\ref{sec:pack} 节；若驱动
黑名单强制回落，HUD 会显示软件渲染器名，本文的数据不是在该模式下取的）。控制方式为
OrbitControls 风格的鼠标交互：左键拖拽旋转、滚轮缩放、
右键拖拽受限平移（默认关闭以保持黑洞居中），
\texttt{H} 键切换侧栏收起后的全屏预览，
\texttt{F} 键进入全屏，\texttt{Escape} 在侧栏收起时还原。
侧栏的开合与画布尺寸彻底解耦，因此无论面板是否可见，
黑洞的阴影圆心始终位于主画面的几何中心；在打包版
可执行文件中实测，侧栏展开/收起两种状态的圆心偏差为
$+0.021$/$+0.042$\,px，均不超过 $0.05$\,px，即约为亚像素探针
标准差（$0.0147$\,px）的 $1.4$ 与 $2.9$ 倍
（第~\ref{sec:viewport} 节、图~\ref{fig:23}）。

与程序一同交付的还有：本文的 \LaTeX{} 源码、由脚本生成的
$19$ 张表与 $25$ 张插图、以及生成它们的全部脚本
（\texttt{paper/tools}）。复现任意一个数字的完整命令清单
见附录~\ref{app:repro}；每张图的可选参数、
每张表的输入文件与输出路径都在该附录中列出。

% --------------------------------------------------------------------------- %
\subsection{展望}

把本文的结论向前推一步，最自然的三个方向是：
\emph{（一）}把背景时空换成 Kerr，用 Bardeen 临界曲线
~\cite{bardeen1973} 与光子壳层的 Lyapunov 指数~\cite{johnson2020,gralla2020}
替换本文的闭式判据，此时第~\ref{sec:discuss-phys} 节列出的
四类效应会同时进入画面，验证体系需要整体重建；
\emph{（二）}把前向累积的体积积分升级为含自吸收与康普顿散射的
完整辐射转移，并用两温或非热电子分布替换本文的局部热平衡假设，
这一步的代价是对每个体积元增加一次频率空间的积分，
预计单帧成本上升 $2$--$4$ 倍（第~\ref{sec:discuss-phys} 节）；
\emph{（三）}把渲染器搬到 WebGPU 计算管线，
用工作组共享内存缓存步长与几何量，把本文的
$n_{\rm max}$ 上限从 $4096$ 推开，从而在
第~\ref{sec:criticalband} 节的对数窄带内获得更多可分辨的
近临界结构。

在这三条之外，本文更希望被读到的是它的\emph{方法论}：
一个实时图形程序完全可以同时具备可证伪的物理、
有残差的数值分析和可一键重算的实验体系。
做到这一点的成本并不高——它需要的不是更快的硬件，
而是在写第一行着色器代码之前，先把闭式解和它的反解路径准备好。
本文的六组验证实验、$19$ 张表与 $25$ 张插图中，
没有一项用于证明程序「看起来像黑洞」；
它们全部用于回答一个更窄、也更有价值的问题：
\emph{当画面与公式不一致时，我们能多快知道是哪一个错了。}
